\chapter{Retrieval-Based Scar Generation}

The previous chapters established a compact geometric characterization of post-infarction
scar. This chapter presents a retrieval-based generation (RBG) framework that leverages
these descriptors to retrieve plausible scar patterns from the available patient cohort.

\section{Problem Formulation}

Given a target specification---which may be a partial clinical profile, a set of desired scar
properties, or an existing patient to be matched---the goal is to produce a plausible scar
polar map that is consistent with the population-level distribution observed in real data.
Formally, given a query vector $\mathbf{q} \in \mathbb{R}^d$ encoding the desired scar
characteristics, the system must return one or more polar maps drawn from real observations
whose descriptor vectors are close to $\mathbf{q}$.

A retrieval-based strategy was chosen over generative alternatives such as VAEs, GANs, or
diffusion models for three main reasons. First, the available cohort is relatively small for
training deep generative models with high-dimensional outputs. Second, every output is, by
construction, a real patient scar observed in clinical practice, which provides a strong
plausibility constraint. Third, retrieval is fully interpretable: the distance to the nearest
neighbour quantifies how well the query is represented in the case base, and the identity and
geometry of the retrieved patient can be inspected directly. The main limitation is that the
method cannot synthesize completely new scar shapes outside the observed cohort; this is
partially addressed through probabilistic retrieval.

\section{Scar Descriptor Space}

The retrieval engine operates in a descriptor space where each patient is represented by a
fixed-length feature vector $\mathbf{v} \in \mathbb{R}^d$. Six candidate descriptor vectors
were compared, designed to evaluate the relative contribution of structured scar features,
AHA segment burdens, dimensionality reduction, and pixel-level polar-map information.

The \textbf{full} vector (39 features) contains all structured descriptors: segment-level scar
burden across the 17 AHA segments, coronary territory fractions, number of involved
territories, circumferential spread, compactness, apex-to-base gradient, number of
fragments, and fragment-level radial and size descriptors. The \textbf{full\_no\_seg} variant
(22 features) removes the 17 AHA segment burdens to test whether territory, morphology,
and fragment descriptors are sufficient without per-segment information. Pruned variants
were derived using SVD contribution analysis: features with contributions below the 25th
percentile were dropped, yielding \textbf{full\_pruned} (29 features) from the full vector and
\textbf{full\_ns\_pruned} (16 features) from the segment-free variant. The
\textbf{pmap\_only} vector (8 features) uses PCA scores from the flattened binary polar
maps, capturing pixel-level spatial variation without manually designed features. The
\textbf{full\_ns+pmap} hybrid (30 features) combines the non-segment structured features
with the polar-map PCA scores.

\section{Similarity Metric and Retrieval Strategy}

Three distance metrics were evaluated on the selected descriptor: standard Euclidean,
Manhattan, and cosine. Manhattan distance on $z$-score normalised features was selected
as the primary metric based on leave-one-out performance.

The retrieval engine follows a $k$-nearest-neighbour strategy. Two modes were implemented:

\textbf{Deterministic retrieval} ($k = 1$) returns the single nearest neighbour:
$i^* = \arg\min_i\, d(\mathbf{q}, \mathbf{v}_i)$.
This mode maximises descriptor-space similarity but always returns the same patient for the
same query.

\textbf{Probabilistic retrieval} ($k = 10$) retains the top 10 neighbours and samples
according to softmax weights:
\begin{equation}
    p_i = \frac{\exp(-d_i/\tau)}{\sum_j \exp(-d_j/\tau)},
    \label{eq:softmax}
\end{equation}
where $\tau > 0$ is a temperature parameter controlling the sharpness of the distribution:
as $\tau \to 0$ all weight concentrates on the nearest neighbour (recovering deterministic
retrieval), while larger $\tau$ spreads weight more uniformly across the top-$k$ candidates,
increasing diversity at the cost of similarity. A value of $\tau = 1.0$ was used throughout:
probabilistic experiments showed that even at this setting the softmax weights were already
strongly concentrated on the nearest neighbour for most patients, leaving very little
probability mass on non-close cases; a smaller $\tau$ would have further reduced diversity
without meaningful gain in retrieval quality. Every output remains a real patient scar,
but repeated calls with the same query can return different plausible cases.

\section{Evaluation Framework}

Retrieval quality was assessed using leave-one-out cross-validation across the full cohort.
For each patient, that patient was temporarily removed from the case base and used as the
query. Three complementary metrics were computed: scar amount error $|\Delta\text{CZ}|$
(absolute difference in core-zone burden between target and retrieved patient), segment
Dice (binarised overlap over the AHA 17-segment model), and polar-map Dice (pixel-level
overlap on the $64 \times 128$ polar grid). All six candidate vectors were evaluated using this
protocol with the same distance metric.

% Figure removed per revision.
